% College Lectures 10-11: primes, primality tests, Fermat/Euler (second versions), factorization, CRT, exponentiation, discrete logs. All numbers verified by script.
\section{Primes, Primality Tests, CRT and Discrete Logs}

\subsection{Counting and checking primes}
\begin{itemize}
  \item Infinitely many primes. $\pi(n)$ $=$ number of primes $\le n$: $\pi(10) = 4$, $\pi(20) = 8$, $\pi(50) = 15$, $\pi(100) = 25$. Bounds (Gauss, Lagrange): $\dfrac{n}{\ln n} < \pi(n) < \dfrac{n}{\ln n - 1.08366}$; for $n = 10^6$: $72{,}382 < \pi(n) < 78{,}543$ (actual 78{,}498).
  \item \textbf{Is $n$ prime?} Test divisibility by every prime $\le \lfloor\sqrt n\rfloor$. $97$: $\lfloor\sqrt{97}\rfloor = 9$, none of 2, 3, 5, 7 divide $\Rightarrow$ prime. $301$: check primes $\le 17$; $7 \mid 301$ ($= 7 \times 43$) $\Rightarrow$ composite.
  \item \textbf{Mersenne} $M_p = 2^p - 1$: prime for $p = 2, 3, 5, 7, 13, 17$ but $M_{11} = 2047 = 23 \times 89$. \textbf{Fermat} $F_n = 2^{2^n} + 1$: $F_0..F_4 = 3, 5, 17, 257, 65537$ are prime; $F_5 = 4294967297 = 641 \times 6700417$.
\end{itemize}

\subsection{Fermat and Euler: second versions and inverses}
\begin{center}\small
\begin{tabular}{@{}lll@{}}
\toprule
 & First version & Second version (no coprime condition)\\
\midrule
Fermat ($p$ prime) & $a^{p-1} \equiv 1 \pmod p$ if $p \nmid a$ & $a^p \equiv a \pmod p$\\
Euler & $a^{\phi(n)} \equiv 1 \pmod n$ if $\gcd(a,n) = 1$ & $a^{k\phi(n)+1} \equiv a \pmod n$ (for $n = pq$)\\
\bottomrule
\end{tabular}
\end{center}
\textbf{Inverses without extended Euclid}: $a^{-1} \equiv a^{p-2} \pmod p$; \;$a^{-1} \equiv a^{\phi(n)-1} \pmod n$.
\begin{center}\small
\begin{tabular}{@{}ll@{\qquad}ll@{}}
\toprule
$8^{-1} \bmod 17 = 8^{15}$ & $= 15$ & $8^{-1} \bmod 77 = 8^{59}$ & $= 29$ \;($\phi = 60$)\\
$5^{-1} \bmod 23 = 5^{21}$ & $= 14$ & $7^{-1} \bmod 15 = 7^{7}$ & $= 13$ \;($\phi = 8$)\\
$60^{-1} \bmod 101 = 60^{99}$ & $= 32$ & $60^{-1} \bmod 187 = 60^{159}$ & $= 53$ \;($\phi = 160$)\\
$22^{-1} \bmod 211 = 22^{209}$ & $= 48$ & $71^{-1} \bmod 100 = 71^{39}$ & $= 31$ \;($\phi = 40$)\\
\bottomrule
\end{tabular}
\end{center}
Examples: $6^{10} \bmod 11 = 1$; \;$3^{12} \bmod 11 = 3^{11} \cdot 3 \equiv 3 \cdot 3 = 9$; \;$6^{24} \bmod 35 = 1$ ($\phi(35) = 24$); \;$20^{62} \bmod 77 = 20^{60+1} \cdot 20 \equiv 20 \cdot 20 = 400 \equiv \mathbf{15}$.

$\phi$ refresher: $\phi(240) = \phi(2^4)\phi(3)\phi(5) = 8 \cdot 2 \cdot 4 = 64$; \;$\phi(49) = 7^2 - 7 = 42$ (\textbf{not} $\phi(7)^2 = 36$: 7 and 7 aren't coprime); $|\mathbb{Z}_{14}^*| = \phi(14) = 6$: $\{1, 3, 5, 9, 11, 13\}$. For $n > 2$, $\phi(n)$ is even.

\subsection{Primality testing}
\textbf{Deterministic}: divisibility test (exponential: $2^{n/2}$ bit operations for an $n$-bit number), AKS (polynomial, $O((\log_2 n)^{12})$). \textbf{Probabilistic}: fast, may wrongly say ``prime''.
\begin{itemize}
  \item \textbf{Fermat test}: if $a^{n-1} \not\equiv 1 \pmod n$, $n$ is composite; if $\equiv 1$, ``probably prime''. Fooled by \textbf{Carmichael numbers}: $2^{560} \equiv 1 \pmod{561}$ but $561 = 3 \cdot 11 \cdot 17$.
  \item \textbf{Square-root test}: if $n$ is prime, the only square roots of 1 mod $n$ are $\pm1$. Mod 7: $\{1, 6\}$; mod 17: $\{1, 16\}$; mod 8: $\{1, 3, 5, 7\}$ (composite); mod 22: only $\{1, 21\}$ even though 22 is composite, so the test is one-way.
\end{itemize}
\begin{keybox}[title=Miller--Rabin test (base $a$)]
Write $n - 1 = m \cdot 2^k$ with $m$ odd. $T = a^m \bmod n$. If $T = \pm1$: \textbf{prime}. Then repeat up to $k - 1$ times: $T \leftarrow T^2 \bmod n$; if $T = +1$: \textbf{composite}; if $T = -1$: \textbf{prime}. If the loop ends: \textbf{composite}.
\end{keybox}
\begin{center}\small
\begin{tabular}{@{}llll@{}}
\toprule
$n$, base & $n - 1 = m \cdot 2^k$ & $T$ sequence & Result\\
\midrule
561, 2 & $35 \cdot 2^4$ & 263, 166, 67, \textbf{1} & composite (correct: Carmichael caught)\\
27, 2 & $13 \cdot 2^1$ & 11 (no loop since $k - 1 = 0$) & composite\\
61, 2 & $15 \cdot 2^2$ & 11, \textbf{60} ($= -1$) & prime\\
4033, 2 & $63 \cdot 2^6$ & 3521, \textbf{4032} ($= -1$) & ``prime'' (wrong: $4033 = 37 \times 109$)\\
4033, 3 & $63 \cdot 2^6$ & 3551, 2443, 3442, 2443, 3442, 2443 & composite\\
\bottomrule
\end{tabular}
\end{center}
\textbf{Recommended practice}: divisibility test by small primes first, then Miller--Rabin with several bases (4033 fools base 2 but not base 3).

\subsection{Factorization}
\textbf{Fundamental theorem}: $n = p_1^{e_1} \cdots p_k^{e_k}$ uniquely. $\gcd = \prod p_i^{\min(a_i, b_i)}$, \;$\operatorname{lcm} = \prod p_i^{\max(a_i, b_i)}$, \;$\gcd \times \operatorname{lcm} = a \times b$.
\begin{itemize}
  \item \textbf{Trial division}: divide by $a = 2, 3, \dots \le \sqrt n$, repeatedly. $1233 = 3^2 \times 137$; $1523357784 = 2^3 \times 3^2 \times 13 \times 37 \times 43987$.
  \item \textbf{Fermat's method}: find $x$ with $x^2 - n = y^2$; then $n = (x + y)(x - y)$. Start $x = \lceil\sqrt n\rceil$ and step up. $n = 5959$: $x = 78$, $w = 125$; $79$: $282$; $80$: $441 = 21^2$ $\Rightarrow 5959 = 101 \times 59$. Fast when the factors are close (for $4033$ it takes 10 steps: $73^2 - 4033 = 36^2$, so $109 \times 37$).
  \item Others (named on slides): Pollard $p - 1$, Pollard rho, quadratic sieve, number field sieve.
\end{itemize}

\subsection{Chinese remainder theorem}
Solves $x \equiv a_1 \pmod{m_1}, \dots, x \equiv a_k \pmod{m_k}$ with \textbf{pairwise coprime} moduli; the solution is unique mod $M$.
\begin{keybox}[title=Four steps]
1.\ $M = m_1 m_2 \cdots m_k$. \;2.\ $M_i = M / m_i$. \;3.\ $M_i^{-1} \bmod m_i$. \;4.\ $x = \left(\sum a_i M_i M_i^{-1}\right) \bmod M$.
\end{keybox}
\begin{center}\small
\begin{tabular}{@{}lllll@{}}
\toprule
System & $M$ & $M_i$ & $M_i^{-1}$ & $x$\\
\midrule
$x \equiv 2 \ (3),\ 3\ (5),\ 2\ (7)$ & 105 & 35, 21, 15 & 2, 1, 1 & $(140 + 63 + 30) \bmod 105 = \mathbf{23}$\\
$x \equiv 3\ (7),\ 3\ (13),\ 0\ (12)$ & 1092 & 156, 84, 91 & 4, 11, 7 & \textbf{276}\\
$S \equiv 1\ (11),\ 12\ (13),\ 2\ (17)$ & 2431 & 221, 187, 143 & 1, 8, 5 & \textbf{155}\\
\bottomrule
\end{tabular}
\end{center}
\textbf{Uses}: secret sharing (the third row: any 3 of the 4 shares $\{1, 12, 2, 3\}$ mod $\{11, 13, 17, 19\}$ recover $S = 155$); arithmetic on big numbers with small moduli: $123 + 334$ via residues mod 99, 98, 97 gives $z \equiv 61, 65, 69$ $\Rightarrow z = 457$; RSA decryption speed-up.

\subsection{Fast exponentiation (square-and-multiply)}
Scan the exponent bits from the least significant: multiply $y$ by the current $a$ when the bit is 1; square $a$ every step. $17^{22} \bmod 21$, $22 = 10110_2$:
\begin{center}\small
\begin{tabular}{@{}ccll@{}}
\toprule
$i$ & $x_i$ & Multiply ($y$, start 1) & Square ($a$, start 17)\\
\midrule
0 & 0 & -- & $a = 17^2 \bmod 21 = 16$\\
1 & 1 & $y = 1 \cdot 16 = 16$ & $a = 16^2 \bmod 21 = 4$\\
2 & 1 & $y = 16 \cdot 4 \bmod 21 = 1$ & $a = 16$\\
3 & 0 & -- & $a = 4$\\
4 & 1 & $y = 1 \cdot 4 = \mathbf{4}$ & (not needed)\\
\bottomrule
\end{tabular}
\end{center}
Cost: about $\log_2 x$ squarings, versus $x$ multiplications naively. (Section~\ref{sec:modexp} has the left-to-right version used in the Neso videos.)

\subsection{Order, primitive roots, discrete logarithms}
\begin{itemize}
  \item \textbf{Order of the group} $|\mathbb{Z}_n^*| = \phi(n)$; e.g.\ $|\mathbb{Z}_{21}^*| = 12$: $\{1, 2, 4, 5, 8, 10, 11, 13, 16, 17, 19, 20\}$.
  \item \textbf{Order of an element} $\operatorname{ord}(a)$: smallest $i$ with $a^i \equiv 1$; it \textbf{divides} $\phi(n)$ (Lagrange). In $\mathbb{Z}_{10}^*$: $\operatorname{ord}(1, 3, 7, 9) = 1, 4, 4, 2$. In $\mathbb{Z}_8^*$: all non-identity elements have order 2, so \textbf{no primitive root}.
  \item \textbf{Primitive root}: $\operatorname{ord}(a) = \phi(n)$. Exist only for $n = 2, 4, p^t, 2p^t$; if they exist there are $\phi(\phi(n))$ of them. $\mathbb{Z}_7^*$: orders $1, 3, 6, 3, 6, 2$ for $a = 1..6$; primitive roots \textbf{3 and 5} ($\phi(6) = 2$ \checkmark).
  \item \textbf{Discrete log}: if $g$ is a primitive root of $p$, every $y \in \mathbb{Z}_p^*$ is $g^x$ for a unique $x \in [1, p-1]$; $x = L_g\, y$. Easy one way (fast exponentiation), believed hard the other way (exhaustive search) $\Rightarrow$ basis of Diffie--Hellman and ElGamal.
\end{itemize}
\begin{center}\small
\begin{tabular}{@{}c|cccccc@{}}
$y$ & 1&2&3&4&5&6\\ \midrule
$L_3 y$ (mod 7) & 6&2&1&4&5&3\\
$L_5 y$ (mod 7) & 6&4&5&2&1&3\\
\end{tabular}
\qquad
\begin{minipage}{0.45\textwidth}\small
$4 \equiv 3^x \pmod 7 \Rightarrow x = L_3 4 = 4$.\\
$6 \equiv 5^x \pmod 7 \Rightarrow x = L_5 6 = 3$.\\
Rules (mod $\phi(n)$): $L_g 1 = 0$, \;$L_g(xy) = L_g x + L_g y$, \;$L_g x^k = k L_g x$.
\end{minipage}
\end{center}

\begin{fixbox}
\begin{itemize}
  \item Slide: ``561 = 33 $\times$ 17'' is true but not the factorisation into primes: $561 = 3 \times 11 \times 17$ (a Carmichael number).
  \item Slide's $\phi(1) = 0$ (Forouzan's convention for ``integers smaller than $n$''); the standard convention is $\phi(1) = 1$. It never matters in RSA.
  \item Euler's second version $a^{k\phi(n)+1} \equiv a$ holds for squarefree $n$ (like $n = pq$ in RSA), not every $n$; e.g.\ it fails for $n = 4$, $a = 2$.
\end{itemize}
\end{fixbox}
